\studynote{1}{Tue 29 Oct 2024 22:35}{Note on SLLN for BRWRE}
\newpage
\subsection{SLLN for BRWRE}

\subsubsection{Preliminary results and assumptions}

Assumptions for BRW:
\begin{enumerate}
	\item The tree is a $k$-ary tree.
\end{enumerate}
Assumptions for the RE:
\begin{enumerate}
	\item Environment is ergodic.
	\item Environment is strongly elliptic.
\end{enumerate}

According to [Comets et al 00], those assumptions on the RE gives the quenched LDP with rate function $I^q_\eta$, and the quenched path-LDP with rate function $I_\eta^{\text{traj}, q}[\phi(\cdot )] \doteq \int _0^1 I_\eta^q \left( \dot{\phi(t)}\right) d t$.


Define the event $B^\varepsilon_v$, have the LDP
\[
	\frac{1}{n} \log P(B^\varepsilon_v) = - I(x^*) - o_{n \to \infty }(1) - o_{\varepsilon \to 0}(1)
.\] 


\subsubsection{The Argument using quenched LDP}



\begin{proposition}
	\label{prop:SLLN-BRWRE}
	Under assumptions stated above, we have that
	\[
		M_n / n \to x^*, \quad \text{ almost surely}
	.\] 
\end{proposition}
\textbf{Upper Bound.} Use same argument as in [Zei20]. 

\textbf{Lower Bound.} We first give the proof for BRW, and then give a modified proof for BRW-RE satisfying above assumptions.

Assume $v \in D_m$ and $S_v = x$. Construct the tree $T^v_\cdot $ as follows. Let $v$ be the root of the tree, $T^v_0 = \{v\} $. Take $T^v_1$ to be the sets of descendants $w$ of $v$ such that $\rho(v, w) = L$ and $S_w > S_v + (x^* - \varepsilon)L$.

More generally, $T^v_j$ consists of all descendants $w$ of $v$ such that there exists $w' \in T^v_{j -1}$ where $\rho(w', w) = L$ and $S_w > S_{w'} + (x^* - \varepsilon) L$.

\textbf{Claim.} The subtree $T^v$ is supercritical for all large $L$.
To see this, calculate
\begin{align*}
	E|T^v_1|
	&= E \sum_{w \in D^{v}_{L}} 1\left( S_w > S_v + (x^* - \varepsilon) L \right)  \\
	&\geq E \sum_{w \in D^{v}_{L}} 1\left( B^{\varepsilon / 3, x^* - \frac{2}{3}\varepsilon}_{v, L}\right)  \\
	&= k^L P\left( B^{\varepsilon / 3, x^* - \frac{2}{3}\varepsilon}_{v, L} \right)  \\
	&= k^L e^{- LI(x^* - 2 \varepsilon / 3) - L o_\varepsilon(1) - L o_L(1)} \\
	&= e^{L(I(x^*) - I(x^* - 2 \varepsilon / 3) - o_\varepsilon(1) - o_L(1))} \\
	&> 1 \qquad \text{for $L$ large}
.\end{align*}

Fix some $L$ such that $T^v$ is supercritical. Now we construct a sequence of subtrees. Assume $k \ge 2$ and the walk is transient to the right or recurrent. Take $v_1$ to be a particle that first hits $1$. Take $v_2$ to be a particle that first hits $2$ and not a descendant of $v_1$. Continue in this fashion.

Now, the sequence of trees $T^{v_i}$ are independent (actually iid up to translation, in the case of BRW with no RE), and each supercritical. By Borel-Cantelli, at least one of them survives. Now follow any $1$-ary subtree of that surviving tree, i.e. path of a particle, we get
\[
	\liminf_{n \to \infty } M_n / n \ge x^* - \varepsilon
.\] 
Take $\varepsilon \to 0$ and we obtain the desired lower bound.

Now we consider the BRW-RE case. The difficulty comes from the fact that $E|T^v_1|$ now depends on the environment in $[S_v - L, S_v + L]$. Even if $E|T^0_1| > 1$, there will be some rare events where $E|T^w_1| \le 1$, so we cannot say that any subtree $T^v$ constructed as above is supercritical.

However, we may introduce the following tempered subtree $\tilde T ^v$ and force it to be supercritical. The construction is as follows:
\begin{itemize}
	\item For any $v \in \tilde T^w$, compute $E|T^v_1|$. We can do so because we are working in quenched environment.
	\item If $E|T^v_1| > 1$, branch as in the construction of $T^v$. That is, the descendants of $v$ in $\tilde T^w$ is $T^v_1$.
	\item If $E|T^v_1|\le 1$, pick a random particle $v'$ in $D^v_L$, and the descendants of $v$ in $\tilde T^w$ is $T^v_1 \cup \{v'\} $.
\end{itemize}

By this construction, $\tilde T^w$ is supercritical. Construct a sequence of independent subtrees as in the BRW case. Then we obtain some subtree $\tilde T^w$ that survives.

Now it suffices to show that, on the event of non-extinction, any infinite trajectory in $\tilde T^w$ has desired limiting velocity. Fix some trajectory $v_j \in T^w_n$ where $v_{j} \in D^{v_{j - 1}}_L$.

Let
\[
	j^* = \sum_{k = 0}^{j - 1} \mathbf{1}\left( E|T_1^{v_k}| > 1 \right) , \qquad j_* = j - j^*
.\] 

Then for all $j > 0$,
\[
	M_{j L} \ge (x^* - \varepsilon) L j^* - L j_*
.\] 
The desired lower bound would follow if we can show
\[
	\lim_{L \to \infty}  \limsup_{j \to \infty } j_* / j  = 0 \quad \text{ a.s.}
.\] 
Using the ergodic theorem we can control the density of bad sites. Define
\[
	B_x := \{[-L + x, L + x] \text{ is bad}\}  := \{E|T^x_1| \le 1\} 
.\] 
Then
\[
	\frac{1}{N} \sum_{n=0}^{N} 1_{B_x} = \frac{1}{N} \sum_{n=0}^{N} 1_{B_0} \circ \theta^n = P(B_0) + o_N(1) = o_L(1) + o_N(1)
.\] 

Thus the density of bad sites is $o_L(1)$.

\underline{This part need to be made rigorous.} The Borel-Cantelli argument gives no information about the starting position of $v_j$. Thus, the event of $v_j$ switching between good and bad sites can be approximated by a Markov chain with two states. Now, we may estimate those probabilities at $0$ :
\begin{itemize}
	\item $P(S_L \,\text{bad} | 0\, \text{bad} ) = o_L(1)$
	\item $P(S_L \,\text{bad} | 0\, \text{good} ) = o_L(1)$
\end{itemize}

Thus, the stationary distribution of the good-bad Markov chain is given by
\[
	\mu(G) = 1 - o_L(1), \qquad \mu(B) = o_L(1)
.\] 
Therefore
\[
	j_* = j \mu(B) + o(j) = j o_L(1) + o(j)
.\] 
This completes the proof.


